% Korean edition of sections/s4_0_parent.tex.
\seclabel
\claim{\secnum{4}{sec:live}\quad 테스트에서 드러난 결함 2개를 고친 뒤, 목표 5개 중 2개는 완전히 확인됐고 G1과 G3는 일부, G4는 미확인이다}

\lead{2026년 11월 24일 실서버 테스트에서 결함 2개가 드러났고, 첫 실서버 실행과 이 테스트 사이의 코드 리뷰에서 다른 결함 7개를 이미 고쳤다.} 테스트는 실제 포트폴리오 하나를 대상으로 09:27\textasciitilde 09:55 KST에 진행했고, 반올림 테스트는 10:26과 10:33에 했다. 테스트 주문은 체결되지 않도록 시장가에서 20\% 떨어진 지정가로 넣었으며, 그날은 새 주문 9건과 정지 7건, 래퍼의 하루 한도 80건 중 16건으로 끝났다. 주문 응답 유실이나 429처럼 실서버에서 일부러 일으킬 수 없는 실패는 오프라인 테스트 31개 중 22개가 다룬다. \ref{sec:wrappertests}절은 시작, 피드, 매수, 매도를 시험하고 \code{plan()}이 합계 139\%인 목표를 받아들였음을 발견한다. \ref{sec:oneshare}절은 서버가 주문 주식 수를 내림함을 보인다. \ref{sec:guidelines}절은 \code{guidelines()}가 잘못된 필드로 섹터 한도를 찾았음을 발견한다. \ref{sec:validity}절은 이 테스트가 보여주지 못하는 것과 남은 질문을 정리한다.
